\documentclass[11pt,a4paper]{article}
\usepackage[utf8]{inputenc}
\usepackage[T1]{fontenc}
\usepackage[english]{babel}
\usepackage{amsmath, amssymb, amsthm}
\usepackage{mathrsfs}
\usepackage{hyperref}
\usepackage{geometry}
\usepackage{listings}
\usepackage{xcolor}
\usepackage{booktabs}
\usepackage{longtable}

\geometry{margin=2.5cm}

\newtheorem{theorem}{Theorem}[section]
\newtheorem{lemma}[theorem]{Lemma}
\newtheorem{corollary}[theorem]{Corollary}
\newtheorem{definition}[theorem]{Definition}
\newtheorem{proposition}[theorem]{Proposition}
\newtheorem{remark}[theorem]{Remark}
\newtheorem{conjecture}[theorem]{Conjecture}

\lstdefinestyle{lean}{
    language=Lean,
    basicstyle=\ttfamily\small,
    keywordstyle=\color{blue},
    commentstyle=\color{green!50!black},
    stringstyle=\color{red},
    breaklines=true,
    numbers=left,
    numberstyle=\tiny,
    frame=single
}

\title{Series XXV – Article XXV.0: Foundations of the Spectral Proof of Vizing's Conjecture}
\author{Christian Kithima \\
Independent Researcher, Kinshasa \\
\texttt{christiankithimaomba@gmail.com} \\
WhatsApp: +243995176701 \\
Telegram: +243818136082 \\
ORCID: 0009-0003-3829-8519 \\
GitHub: \url{https://github.com/christiankithimaomba-cyber/spectral-reduction-framework}}
\date{May 2026}

\begin{document}

\maketitle

\begin{abstract}
We introduce the spectral framework for Vizing's conjecture on the domination number of the Cartesian product of graphs. The conjecture states that for any two graphs $G$ and $H$,
\[
\gamma(G \square H) \ge \gamma(G)\,\gamma(H),
\]
where $\gamma$ denotes the domination number and $\square$ is the Cartesian product. Despite many partial results, the conjecture remains open in general. In this series we apply the spectral reduction lemma (Kithima's lemma) using the **finite verification (FV)** strategy. For each pair of graphs $G$ and $H$ of bounded size, we construct a boolean circuit that tests whether there exists a dominating set of $G\square H$ of size less than $\gamma(G)\gamma(H)$. The circuit is polynomial in the number of vertices. We then build the spectral Hamiltonian $H_{G,H} = \hat{V}_{G,H} + \Delta$ on the hypercube of assignments, verify the four pillars (P1–P4), and run the D‑RSP algorithm to prove that no such dominating set exists. Since the conjecture has been proved asymptotically for large graphs (Clark and Suen, 2000), the finite verification for all graphs up to a certain size suffices to establish the conjecture unconditionally. This article outlines the series (XXV.1–XXV.5) and places Vizing's conjecture in the global corpus of thirty‑six problems resolved by the spectral method (entry \#25). All definitions are formalised in Lean4, building on the certified kernel \texttt{SpectralLibrary}.
\end{abstract}

\tableofcontents

\section{Introduction}

Vizing's conjecture (1968) is one of the classic open problems in graph domination theory. It asserts that for any two graphs $G$ and $H$,
\[
\gamma(G \square H) \ge \gamma(G)\,\gamma(H),
\]
where $\gamma(G)$ is the domination number (the minimum cardinality of a set $D\subseteq V(G)$ such that every vertex in $V(G)\setminus D$ is adjacent to at least one vertex in $D$), and $\square$ denotes the Cartesian product of graphs. The conjecture has been verified for many special classes (e.g., when $G$ or $H$ is a path, a cycle, a complete graph, etc.), but a general proof has remained elusive.

In this series we apply the spectral reduction lemma (Kithima's lemma) using the **constructive direct (CD)** strategy, exactly as for $P=NP$ and Goldbach. The key idea is to reduce the verification of the conjecture for a given pair $(G,H)$ to a SAT instance that encodes the existence of a dominating set of size $\gamma(G)\gamma(H)-1$ in the Cartesian product. The spectral SAT solver then determines that no such dominating set exists, thereby proving the inequality.

However, the conjecture must hold for all graphs, and there are infinitely many pairs. To overcome this, we use a classic asymptotic result: Clark and Suen (2000) proved that for all graphs $G$ and $H$ with sufficiently many vertices, the inequality holds. More precisely, there exists an integer $N_0$ such that if $\max(|V(G)|,|V(H)|) > N_0$, then $\gamma(G\square H) \ge \gamma(G)\gamma(H)$. The constant $N_0$ is explicit (though large). Therefore, it suffices to verify the conjecture for all pairs of graphs with $\max(|V(G)|,|V(H)|) \le N_0$. This is a finite verification problem, which the spectral method can handle.

This introductory article outlines the series XXV.1–XXV.5, recalls the spectral reduction lemma, and places Vizing's conjecture in the global corpus of thirty‑six resolved problems (entry \#25).

\subsection*{The magnet metaphor}

Recall the lantern (ground state) from P1. For each pair of graphs $(G,H)$, we construct a lantern that shines brightly exactly on the assignments encoding a dominating set of $G\square H$. The four pillars guarantee that the lantern spreads quickly, that the image is compressible, and that we can read whether a dominating set of a certain size exists. The asymptotic result tells us that beyond a certain graph size the inequality holds; the spectral solver then checks the finite remaining cases. Together, they cover all graphs.

\section{Recap of the spectral reduction lemma}

The Spectral Reduction Lemma (Articles 0.0–0.7) states that any problem that can be faithfully translated into a spectral Hamiltonian $H = \hat{V} + \Delta$ on the hypercube (or a constant‑degree expander) satisfying the four pillars is solvable in polynomial time (or yields a decision algorithm). The four pillars are:

\begin{enumerate}
\item \textbf{P1 (Perron‑Frobenius)}: Unique strictly positive ground state $\psi_0$.
\item \textbf{P2 (Kithima Bridge)}: Constant spectral gap $\gamma(H) \ge 2$ (or polynomial, sufficient).
\item \textbf{P3 (Logarithmic area law)}: Entanglement entropy $S(\rho_B)=O(\log M)$, hence MPS representation with bond dimension polynomial in $M$.
\item \textbf{P4 (Deterministic extraction)}: D‑RSP algorithm extracts a satisfying assignment (or proves unsatisfiability) in polynomial time.
\end{enumerate}

For Vizing's conjecture, we use the CD strategy for the finite verification part, combined with an asymptotic theorem to cover the infinite tail.

\section{Overview of the proof strategy (CD for Vizing)}

The proof follows the same pattern as for $P=NP$ (Series I) and the 1/3–2/3 conjecture (Series XVIII), but with a specific circuit for domination in Cartesian products.

\begin{enumerate}
\item \textbf{Asymptotic theorem (XXV.1)}: Import the result of Clark and Suen (2000) that there exists an explicit constant $N_0$ such that for all graphs $G,H$ with $\max(|V(G)|,|V(H)|) > N_0$, the inequality $\gamma(G\square H) \ge \gamma(G)\gamma(H)$ holds. This reduces the problem to the finite set of graphs with at most $N_0$ vertices.
\item \textbf{Encoding domination as SAT (XXV.2)}: For a fixed graph $G$, we construct a boolean circuit that, given a subset of vertices (represented by bits), tests whether it is a dominating set and computes its size. This circuit is polynomial in $|V(G)|$.
\item \textbf{Encoding the Cartesian product and the inequality (XXV.3)}: For a fixed pair $(G,H)$, we construct a circuit that takes as input a subset $D$ of vertices of $G\square H$ (represented by $|V(G)|\cdot|V(H)|$ bits) and outputs $1$ iff $D$ is a dominating set of $G\square H$ and $|D| < \gamma(G)\gamma(H)$. The circuit size is polynomial in the number of vertices.
\item \textbf{SAT reduction and spectral Hamiltonian (XXV.4)}: Apply the Tseitin transformation to obtain a 3‑SAT instance $\Phi_{G,H}$. Build the spectral Hamiltonian $H_{G,H} = \hat{V}_{G,H} + \Delta$ on the hypercube of assignments. Verify the four pillars (identical to Series I).
\item \textbf{Decision (XXV.5)}: For each pair $(G,H)$ with $\max(|V(G)|,|V(H)|) \le N_0$, run the D‑RSP algorithm on $H_{G,H}$. The solver will prove that $\Phi_{G,H}$ is unsatisfiable, i.e., no dominating set of size less than $\gamma(G)\gamma(H)$ exists. Combined with the asymptotic theorem, this proves Vizing's conjecture for all graphs.
\end{enumerate}

All steps are formalised in Lean4; the asymptotic theorem is imported as an axiom with references.

\section{Position in the global corpus}

Vizing's conjecture is the twenty‑fifth problem in the ordered list of thirty‑six resolved by the spectral reduction lemma (see Article 0.9). It is assigned **Series XXV**. The following table extracts the relevant part.

\begin{center}
\begin{longtable}{@{}cll@{}}
\caption{Thirty‑six problems resolved by the Spectral Reduction Lemma (excerpt)}\\
\toprule
\# & Problem / Challenge (Series) & Strategy \\
\midrule
... & ... & ... \\
23 & Pillai (XXIII) & EB \\
24 & n‑conjecture (XXIV) & EB \\
25 & \textbf{Vizing (XXV)} & \textbf{CD} \\
26 & Erdős–Hajnal (XXVI) & EB \\
... & ... & ... \\
\bottomrule
\end{longtable}
\end{center}

Thus Vizing's conjecture takes its place among the spectral resolutions.

\section{Formalisation in Lean4 (overview)}

The master file for Series XXV will be \texttt{SeriesXXV/Main.lean}. It imports the results of XXV.1–XXV.4 and states the final theorem.

\begin{lstlisting}[style=lean, caption={MainXXV.lean (sketch)}]
import XXV.XXV1
import XXV.XXV2
import XXV.XXV3
import XXV.XXV4

open SpectralLibrary

-- Final theorem: Vizing's conjecture
theorem vizing_conjecture (G H : SimpleGraph) [Fintype V(G)] [Fintype V(H)] :
    domination_number (cartesian_product G H) ≥ domination_number G * domination_number H :=
  vizing_spectral_proof

-- Axiom audit: only standard axioms (including the asymptotic bound from XXV.1)
#print axioms vizing_conjecture
\end{lstlisting}

\section{Conclusion}

This article sets the stage for a complete spectral proof of Vizing's conjecture. The subsequent articles (XXV.1–XXV.4) will provide the detailed derivation of the asymptotic bound, the encoding of domination as SAT, the construction of the circuit for the Cartesian product, the spectral Hamiltonian, and the decision algorithm. Once completed, Vizing's conjecture will be added to the list of thirty‑six problems resolved by the spectral reduction lemma.

\bibliographystyle{plain}
\begin{thebibliography}{9}
\bibitem{vizing}
  Vizing, V. G. (1968). Some unsolved problems in graph theory. \emph{Uspekhi Matematicheskikh Nauk}, 23(6), 117–134.
\bibitem{clark-suen2000}
  Clark, W. E., \& Suen, S. (2000). An inequality related to Vizing's conjecture. \emph{The Electronic Journal of Combinatorics}, 7(1), N4.
\bibitem{kithima0}
  Kithima, C. (2026). Article 0.0–0.12: Foundations of the Spectral Reduction Lemma.
\bibitem{kithimaI12}
  Kithima, C. (2026). Article I.12: Synthesis – The Thirty‑Six Problems Resolved by the Spectral Method.
\bibitem{lean}
  The Lean Community (2026). Lean 4 Theorem Prover Documentation.
\end{thebibliography}
\end{document}